%============================================================
% CH2 · LESSON 3 — THE DFS TREE AND THE FOUR EDGE TYPES
%============================================================
\subsection{The DFS Tree and the Four Edge Types}
\label{ssec:ch2-dfs-tree}

Every vertex except the start vertices is entered through exactly one edge: the edge $u \to v$ that DFS was examining when it found $v$ white. These edges form the \textbf{DFS tree} (one tree per start vertex of the outer loop --- together a \textbf{DFS forest}). The other edges are not wasted: each of them was examined once and found a vertex that was gray or black, and \emph{which} colour it found says a lot about the graph.

\begin{definitionbox}[The four edge types (directed graphs)]
When DFS, sitting at $u$, examines the edge $u \to v$:
\begin{center}
\small
\renewcommand{\arraystretch}{1.3}
\begin{tabular}{@{}llll@{}}
\toprule
\textbf{colour of $v$} & \textbf{extra test} & \textbf{type} & \textbf{meaning in the DFS forest} \\
\midrule
white & --- & \textbf{tree} edge & $v$ becomes a child of $u$ \\
gray & --- & \textbf{back} edge & $v$ is an ancestor of $u$ (or $v = u$, a self-loop) \\
black & $\tin{u} < \tin{v}$ & \textbf{forward} edge & $v$ is a descendant of $u$, but not its child \\
black & $\tin{v} < \tin{u}$ & \textbf{cross} edge & neither is an ancestor of the other \\
\bottomrule
\end{tabular}
\end{center}
\end{definitionbox}

Why these meanings? A gray $v$ is on the stack, i.e.\ on the path from the root to $u$: an ancestor. A black $v$ with $\tin{v} > \tin{u}$ was entered while $u$ was gray, so it is a descendant (parenthesis theorem), and it is not a child since $u \to v$ is not its tree edge. A black $v$ with $\tin{v} < \tin{u}$ was entered before $u$ and is already finished, so its interval ends before $u$'s begins: the two are unrelated.

\begin{diagrambox}
\begin{tikzpicture}[x=1.15cm, y=1cm]
    % the DFS forest of D, drawn as rooted trees
    \node[gvsrc] (t1) at (2,0) {1};
    \node[gv] (t2) at (0.8,-1.2) {2};
    \node[gv] (t5) at (3.2,-1.2) {5};
    \node[gv] (t3) at (0.1,-2.4) {3};
    \node[gv] (t4) at (1.5,-2.4) {4};
    \node[gv] (t6) at (3.2,-2.4) {6};
    \node[gv] (t7) at (3.2,-3.6) {7};
    \node[gvsrc] (t8) at (5.2,0) {8};
    \foreach \a/\b in {1/2, 1/5, 2/3, 2/4, 5/6, 6/7} \draw[gtree, -{Stealth[length=2.4mm]}] (t\a) -- (t\b);
    \draw[gback] (t3) to[bend left=40] (t1);
    \draw[gback] (t7) to[bend right=50] (t6);
    \draw[gfwd] (t1.east) .. controls (4.5,-0.3) and (4.4,-2.0) .. (t6.east);
    \draw[gcross] (t4) -- (t3);
    \draw[gcross] (t5) to[bend left=15] (t4);
    \draw[gcross] (t8) to[bend left=30] (t7);
    \foreach \v/\a/\b/\p in {1/0/13/left, 2/1/6/left, 3/2/3/below, 4/4/5/below, 5/7/12/left, 6/8/11/left, 7/9/10/left, 8/14/15/right}
        \node[gdist, text=darkgray, \p=0.5mm of t\v] {\tiny[\a,\b]};
    % legend
    \node[anchor=west, font=\scriptsize, align=left] at (6.8,-0.6) {%
        \tikz[baseline=-0.5ex]\draw[gtree, -{Stealth[length=2mm]}] (0,0) -- (0.7,0);\ tree: $1{\to}2$, $2{\to}3$, $2{\to}4$, $1{\to}5$, $5{\to}6$, $6{\to}7$\\[3pt]
        \tikz[baseline=-0.5ex]\draw[gback] (0,0) -- (0.7,0);\ back: $3{\to}1$, $7{\to}6$\\[3pt]
        \tikz[baseline=-0.5ex]\draw[gfwd] (0,0) -- (0.7,0);\ forward: $1{\to}6$\\[3pt]
        \tikz[baseline=-0.5ex]\draw[gcross] (0,0) -- (0.7,0);\ cross: $4{\to}3$, $5{\to}4$, $8{\to}7$};
    \node[anchor=west, font=\scriptsize\itshape, text=darkgray, text width=6.3cm] at (6.8,-2.7)
        {Children are drawn left to right in the order DFS entered them. Every cross edge points to the \textbf{left} (to an earlier, finished part of the forest) --- never to the right.};
\end{tikzpicture}
\figcap{The DFS forest of $D$ with every one of its 12 edges classified (the program below prints exactly this). Green border $=$ root of a DFS tree; small numbers $= [\tin{}, \tout{}]$.}
\end{diagrambox}

\begin{insightbox}[Cross edges only point back in time]
An edge $u \to v$ with $\tin{u} < \tin{v}$ can never be a cross edge: when DFS examines it, $v$ is either white (tree edge), or it was entered after $u$ and hence is $u$'s descendant (forward edge). So a cross edge always goes from a later vertex to an earlier, already finished one --- in pictures, from right to left. This asymmetry is used in Chapter~5 (topological order) and Chapter~19 (strongly connected components).
\end{insightbox}

\subsubsection*{Undirected graphs: only tree and back edges}

In an undirected graph each edge $\{u, v\}$ is stored twice and DFS examines it from both ends. Classify an edge by the \emph{first} time it is examined.

\begin{theorembox}[No forward or cross edges in undirected graphs]
In a DFS of an undirected graph every edge is a tree edge or a back edge (it joins a vertex to one of its ancestors).
\end{theorembox}

\textbf{Proof.} Let $\{u, v\}$ be an edge with $\tin{u} < \tin{v}$. While $u$ is gray, DFS scans all of $u$'s neighbours, including $v$; if $v$ is still white then, $\{u,v\}$ becomes a tree edge. Otherwise $v$ was entered while $u$ was gray, so $v$ is a descendant of $u$. Then $v$ is entered and scanned before $u$ is left, and the edge is first examined from $v$'s side, where it leads to the gray ancestor $u$: a back edge. A cross edge would need two vertices with disjoint intervals joined by an edge --- but the earlier one would have pulled the later one into its subtree. \hfill$\blacksquare$

\begin{diagrambox}
\begin{tikzpicture}[x=1.15cm, y=1cm]
    \begin{scope}
        \pic {g2-G3};
        \node[font=\scriptsize\bfseries] at (2.15,-1.8) {$G_3$ (Chapter 1)};
    \end{scope}
    \begin{scope}[xshift=7.2cm, yshift=1.0cm]
        \node[gvsrc] (u1) at (0,0) {1};
        \node[gv] (u2) at (0,-1) {2};
        \node[gv] (u3) at (-0.9,-2) {3};
        \node[gv] (u5) at (0.9,-2) {5};
        \node[gv] (u4) at (0.3,-3) {4};
        \node[gv] (u6) at (1.5,-3) {6};
        \node[gvsrc] (u7) at (2.6,0) {7};
        \foreach \a/\b in {1/2, 2/3, 2/5, 5/4, 5/6} \draw[gtree] (u\a) -- (u\b);
        \draw[gback, -] (u4) to[bend left=55] (u1);
        \node[anchor=west, font=\scriptsize, text width=4.4cm] at (2.3,-1.7) {DFS from 1 (then from 7): five tree edges and one back edge $4$--$1$, which joins 4 to its ancestor 1. No edge joins two different branches.};
    \end{scope}
\end{tikzpicture}
\figcap{The undirected graph $G_3$ of Lesson~\ref{ssec:ch1-first-traversal} and its DFS forest. The edge $4$--$5$ is a tree edge (5 entered 4); the edge $1$--$4$ is first examined from 4, where 1 is gray: a back edge.}
\end{diagrambox}

\begin{warningbox}[The edge back to the parent is not a back edge]
From $v$'s side, the tree edge $\{p, v\}$ to the parent $p$ leads to a gray vertex too. It must be \textbf{skipped}, otherwise every tree edge looks like a cycle of length 2. Skip it by remembering the \emph{edge id} through which $v$ was entered, not the parent \emph{vertex}: with two parallel edges between $p$ and $v$, the second one really is a back edge (a cycle of length 2), and ``skip the parent vertex'' would miss it. Lesson~\ref{ssec:ch2-cycles} builds cycle detection on exactly this rule.
\end{warningbox}

\subsubsection*{Code: classify every edge}

\begin{lstlisting}[style=cpp]
#include <bits/stdc++.h>
using namespace std;

// Classify every edge of a DIRECTED graph as tree / back / forward / cross.
vector<vector<int>> adj;
vector<int> color, tin;
int timer = 0;
map<string, int> cnt;

void dfs(int u) {
    color[u] = 1; tin[u] = timer++;
    for (int v : adj[u]) {
        string type;
        if (color[v] == 0) type = "tree";
        else if (color[v] == 1) type = "back";           // v is on the stack
        else type = (tin[u] < tin[v]) ? "forward" : "cross";
        cout << u << "->" << v << " " << type << "\n";
        cnt[type]++;
        if (color[v] == 0) dfs(v);
    }
    color[u] = 2; timer++;                               // the exit tick (tout) is not needed here
}

int main() {
    int n, m;
    cin >> n >> m;
    adj.assign(n + 1, vector<int>());
    for (int i = 0; i < m; i++) {
        int a, b;
        cin >> a >> b;
        adj[a].push_back(b);
    }
    color.assign(n + 1, 0); tin.assign(n + 1, -1);
    for (int s = 1; s <= n; s++)
        if (color[s] == 0) dfs(s);
    for (auto& p : cnt) cout << p.first << ": " << p.second << "\n";
}
\end{lstlisting}

\begin{lstlisting}[style=trace]
input : D
output: 1->2 tree      2->3 tree      3->1 back      2->4 tree
        4->3 cross     1->5 tree      5->4 cross     5->6 tree
        6->7 tree      7->6 back      1->6 forward   8->7 cross
        back: 2   cross: 3   forward: 1   tree: 6        (printed one per line)
\end{lstlisting}

For undirected graphs the same program needs edge ids, to skip the parent edge and to report each edge once:

\begin{lstlisting}[style=cpp]
#include <bits/stdc++.h>
using namespace std;

// Undirected graph: every edge is a tree edge or a back edge. adj stores (neighbour, edge id).
vector<vector<pair<int, int>>> adj;
vector<int> color, parentEdge;
vector<bool> done;                                       // edge already classified?

void dfs(int u) {
    color[u] = 1;
    for (auto e : adj[u]) {
        int v = e.first, id = e.second;
        if (id == parentEdge[u] || done[id]) continue;   // skip the edge we came in by
        done[id] = true;
        if (color[v] == 0) {
            cout << u << "-" << v << " tree\n";
            parentEdge[v] = id;
            dfs(v);
        } else {
            cout << u << "-" << v << " back\n";          // v is gray: an ancestor of u
        }
    }
    color[u] = 2;
}

int main() {
    int n, m;
    cin >> n >> m;
    adj.assign(n + 1, {});
    for (int i = 0; i < m; i++) {
        int a, b;
        cin >> a >> b;
        adj[a].push_back({b, i});
        adj[b].push_back({a, i});
    }
    color.assign(n + 1, 0); parentEdge.assign(n + 1, -1); done.assign(m, false);
    for (int s = 1; s <= n; s++)
        if (color[s] == 0) dfs(s);
}
\end{lstlisting}

\begin{lstlisting}[style=trace]
input : G3 (7 6 / 1 2 / 1 4 / 2 3 / 2 5 / 4 5 / 5 6)
output: 1-2 tree   2-3 tree   2-5 tree   5-4 tree   4-1 back   5-6 tree
\end{lstlisting}

\begin{comparisonbox}[What each edge type is used for]
\small
\begin{tabularx}{\linewidth}{@{}>{\bfseries}l X@{}}
tree & the DFS tree itself: parents, subtrees, depths (Patterns~2 and~3) \\
back & \textbf{cycles}: a back edge $u \to v$ plus the tree path $v \leadsto u$ is a cycle (Lesson~\ref{ssec:ch2-cycles}, Pattern~1); in undirected graphs, low-links, bridges and cut vertices (Chapter~18) \\
forward & harmless for cycles (the tree path already goes there); matters only in algorithms that care about multiple paths \\
cross & harmless for cycles; important for strongly connected components (Chapter~19) \\
\end{tabularx}
\end{comparisonbox}

\subsubsection*{Check your understanding}

\begin{enumerate}[leftmargin=*]
    \item A DFS of a connected undirected graph with $n$ vertices and $m$ edges: how many tree edges and how many back edges does it find?
    \item Can a forward edge appear in a DFS of an undirected graph? Can a self-loop $u \to u$ in a directed graph be anything other than a back edge?
    \item In $D$, the edge $8 \to 7$ is a cross edge. Restart the outer loop at $s = 8$ first (then $1, 2, \dots$). What is the type of $8 \to 7$ now, and of $6 \to 7$?
    \item Why does the classification of a directed edge need $\tin{}$ only for black targets?
    \item A directed graph has a DFS with no back edges. What can you say about it?
\end{enumerate}

\begin{solutionbox}[Answers]
\begin{enumerate}[leftmargin=*]
    \item $n - 1$ tree edges (one per vertex except the root) and $m - n + 1$ back edges --- every other edge.
    \item No (theorem above). A self-loop finds $u$ itself, which is gray: always a back edge, i.e.\ a cycle of length 1.
    \item DFS from 8 enters 7 via $8 \to 7$ (\textbf{tree}) and 6 via $7 \to 6$ (tree); from 6 the edge $6 \to 7$ finds 7 gray: a \textbf{back} edge. The later search from 1 finds 6 black and entered earlier through both $5 \to 6$ and $1 \to 6$: two cross edges (so $1 \to 6$ changes from forward to cross). The cycle $6 \rightleftarrows 7$ is still detected, through a different edge. (Output of the trace program with the start order $8, 1, 2, \dots$.)
    \item White means tree and gray means back, whatever the times. Only for a black target are two answers possible, and the order of entry decides between descendant (forward) and unrelated (cross).
    \item It is acyclic (a DAG): the next lesson proves that a directed graph has a cycle iff some DFS finds a back edge --- and then \emph{every} DFS finds one.
\end{enumerate}
\end{solutionbox}
